\documentclass[11pt]{article}
\usepackage[T1]{fontenc}
\usepackage{textcomp}
\usepackage[margin=0.75in]{geometry}
\usepackage{amsmath,amssymb,amsthm}
\usepackage{array,booktabs}
\usepackage{hyperref}
\setlength{\parindent}{0pt}
\newtheorem{theorem}{Theorem}
\theoremstyle{definition}
\newtheorem{definition}{Definition}
\title{Flow Proof Artifact: prop-03-difference-of-equimultiples}
\author{Generated by \texttt{flow doc proof}}
\date{Flow Proof Book}
\begin{document}
\maketitle
If a first magnitude is the same multiple of a second that a third is of a fourth, the difference of the first and third is the same multiple of the difference of the second and fourth.\\[0.5em]
\textbf{Source.} euclid — Elements, Book V, Proposition 3\\[0.5em]
\bigskip
\noindent{\large \textbf{Derived fact 1.} \textit{If a first magnitude is the same multiple of a second that a third is of a fourth, the difference of the first and third is the same multiple of the difference of the second and fourth} \hfill the Euclidean plane \cdot Euclid Book V \cdot Proposition 3: the difference of equimultiples is an equimultiple of the difference}\par
\smallskip\noindent\textit{Coordinate: the Euclidean plane \cdot Euclid Book V \cdot Proposition 3: the difference of equimultiples is an equimultiple of the difference}\par
\smallskip\noindent\textit{Source: euclid — Elements, Book V, Proposition 3}\par
\smallskip\noindent\textit{Built on: proposition 2: the sum of equimultiples is an equimultiple of the sum, for Euclid Book V on the Euclidean plane}\par
\medskip
\noindent\textbf{Goal.} If a first magnitude is the same multiple of a second that a third is of a fourth, the difference of the first and third is the same multiple of the difference of the second and fourth\par
\noindent$\displaystyle \text{difference equimultiple preserved}$\par
\medskip
\noindent
\begin{tabular}{@{} >{\raggedright\arraybackslash}p{0.06\textwidth} >{\raggedright\arraybackslash}p{0.47\textwidth} >{\raggedleft\arraybackslash}p{0.41\textwidth} @{}}
\textbf{\#} & \textbf{Proof} & \textbf{Mathematics} \\
\midrule
\textbf{1.} & We prove that proposition 3: the difference of equimultiples is an equimultiple of the difference for Euclid Book V the Euclidean plane. & \\[0.45em]
\textbf{2.} & Let a = first. & \\[0.45em]
\textbf{3.} & Let b = second. & \\[0.45em]
\textbf{4.} & Let c = third. & \\[0.45em]
\textbf{5.} & Let d = fourth. & \\[0.45em]
\textbf{6.} & From step 2, step 3, step 4, and step 5, we can deduce that equimultiple(a - c, b - d). & $\displaystyle equimultiple(a - c, b - d)$ \\[0.45em]
\textbf{7.} & We invoke the derived fact governing Euclid Book V the Euclidean plane: proposition 2: the sum of equimultiples is an equimultiple of the sum, for Euclid Book V on the Euclidean plane. & \\[0.45em]
\textbf{8.} & From step 7, this implies difference equimultiple preserved. Hence proven. & $\displaystyle \text{difference equimultiple preserved}$ \\[0.45em]
\bottomrule
\end{tabular}
\label{thm:Geometry-Euclid Book V-proposition_3_the_difference_of_equimultiples_is_an_equimultiple_of_the_difference}
\par\medskip\hrule
\end{document}
